\subsection{已有结果与需要解释的机制}
\label{sec:evidence}
研究首先由实机现象驱动。表\ref{tab:preliminary}记录用户目前提供的近似成功率，
而非本稿重新汇总的最终统计结果。
\begin{table}[t]
 \centering
 \caption{初步实机结果；次数、视角及配对条件待逐项核对}
 \label{tab:preliminary}
 \begin{tabular}{lr}
 \toprule
 方法描述 & 近似成功率\\
 \midrule
 基础ACT & 50\%\\
 ACT + MLP残差 & 60\%\\
 ACT + GRU残差 & 70\%\\
 语义输入策略（无QToken） & 70\%\\
 双视角语义 + QToken & 80\%\\
 语义 + QToken + QTokenGRU & 87\%\\
 \bottomrule
 \end{tabular}
\end{table}
用户已确认各组使用相同的策略示范数据与训练轮次，
因此该结果是在一致的数据与训练预算下比较不同设计，而不是更多动作示范带来的收益。
这里语义输入不是语义替代全部RGB，具体视觉配置以checkpoint为准。
不能假设各行具有相同分母、训练随机性或相同基础策略，
也不能将约10个百分点的差异直接解释为显著的因果效应。

这些结果提出三项可检验解释。
\emph{表示解释：}语义预训练将对象区域的先验显式提供给策略，
几何监督进一步减少目标关系被编码丢失的风险。
应检查这种收益是否集中于拨空减少、有效接触和完整入区，
并排除单纯增加视角或token容量。
\emph{反馈解释：}当前帧残差能根据最新证据修正当前计划，
因此无需每次重选完整动作块。
应比较相同基础checkpoint下的无残差、MLP残差与频繁ACT查询。
\emph{历史解释：}GRU可借助过去的观测和动作上下文区分当前画面中的运动歧义。
应与历史置乱、仅最新帧及同历史的非循环模型比较，
而不把GRU对MLP的所有收益都归于记忆架构。

从约80\%到87\%还提示几何表示与纠偏可能互补，
但证明完整系统更好不等于证明模块之间存在协同效应。
后续优先以相同平台、视角和预算下的
RGB/语义QToken与无残差/GRU组成匹配的 $2\times2$ 对照，
检验收益是否仍然存在；未完成这一控制前不报告加性分解或协同结论。
理论的职责是把上述解释转为明确条件与反证实验，
不是用已经观察到的收益反推所有假设必然成立。
